%=======================================================================
% Related Work, reframed for IEEE Internet of Things Journal (IoT-J).
% Drop-in replacement for Sec. II of ../main.tex: %=======================================================================
% Related Work, reframed for IEEE Internet of Things Journal (IoT-J).
% Drop-in replacement for Sec. II of ../main.tex: %=======================================================================
% Related Work, reframed for IEEE Internet of Things Journal (IoT-J).
% Drop-in replacement for Sec. II of ../main.tex: %=======================================================================
% Related Work, reframed for IEEE Internet of Things Journal (IoT-J).
% Drop-in replacement for Sec. II of ../main.tex: \input{iotj/related_work_iotj}
% Needs ../tables/gap_iotj.tex (capability table) and
% ../tables/gap_fill_iotj.tex (gap -> what SCooP-iT provides -> evidence).
% New bib entries are appended to ../references.bib under "IoT-J additions".
% IoT-J readers come from wireless/IoT systems, not robotics, so the order
% is: wireless testbeds -> CSI sensing -> networked robots -> cooperative
% perception -> indoor ground truth -> trace-driven emulation.
%=======================================================================
\section{Related Work}
\label{sec:related}

\subsection{Wireless and IoT Experimentation Testbeds}
Open testbeds made wireless and IoT experiments reproducible at scale.
FIT IoT-LAB \cite{adjih2015iotlab} offers thousands of low-power nodes
and over a hundred mobile robots across several indoor sites; the PAWR
platforms POWDER \cite{breen2021powder} and COSMOS
\cite{raychaudhuri2020cosmos} provide city-scale software-defined radio,
massive MIMO, and edge computing; Colosseum
\cite{bonati2021colosseum} emulates the channel in hardware for
hundreds of radios. These testbeds treat the radio as the object of
study: they measure the link precisely, but carry no perception
sensors, no per-node pose reference, and no task whose output depends
on what crosses the link. Their mobile nodes are carriers of radios,
not agents whose position and observations are scored.

\subsection{Wi-Fi Channel State Information as a Measurement}
Commodity CSI extraction, first on the Intel~5300
\cite{halperin2011csitool}, then on Broadcom chips with Nexmon
\cite{nexmon} and across several NIC families with PicoScenes
\cite{jiang2022picoscenes}, turned Wi-Fi into a sensor
\cite{ma2019wifisensing}. Channel-sounding and sensing datasets such as
DICHASUS \cite{dichasus}, ESPARGOS \cite{espargos}, MM-Fi \cite{mmfi},
and DeepSense~6G \cite{deepsense6g} pair CSI or beam power with
positions or other modalities, and ray-traced digital twins
\cite{hoydis2023sionnart,alkhateeb2023digitaltwin,bostontwin,deepmimo}
and learned radio maps \cite{levie2021radiounet} predict the channel
from geometry. In this line, CSI serves sensing or localization of a
target, and the geometry is either surveyed or simulated. None records
CSI on the links that a robot team uses to exchange perception data,
on the same clock as the robots' sensors, and none checks a measured
map against the measured channel as two independent references.

\subsection{Networked and Communication-Aware Multi-Robot Systems}
The Internet of Robotic Things \cite{simoens2018iort} treats robots as
mobile IoT nodes that sense, act, and share data over the same network
as other devices. Communication-aware robotics plans motion around the
channel \cite{muralidharan2021commaware}: robots map signal strength
online \cite{fink2010radiomap} or exploit the directionality of the
received signal to keep a team connected \cite{gil2015adaptive}.
ROS-NetSim \cite{calvofullana2021rosnetsim} couples a robot simulator
with a network simulator; CHORD and ACHORD
\cite{ginting2021chord,saboia2022achord} and MOCHA
\cite{cladera2024mocha} share data over intermittent real links in
subterranean and field deployments; \cite{blumenkamp2022framework}
provides a framework for real-robot wireless experiments; and
\cite{zhang2024ros2middleware} compares DDS, Zenoh, and MQTT for ROS~2
over Ethernet, Wi-Fi, and 4G. Multi-robot testbeds
\cite{wilson2021robotarium,paull2017duckietown,kloock2021cpmlab,testa2021choirbot}
give precise positions from motion capture or overhead cameras but
treat the network as ideal. These works model, survive, or benchmark
the link; none logs it per pair, at the rate of the sensors and on
their clock, in a form that can be replayed into another method, and
none scores perception against an independent reference.

\subsection{Cooperative Perception Under Imperfect Communication}
Cooperative perception fuses raw points \cite{chen2019cooper},
features \cite{chen2019fcooper,wang2020v2vnet}, or selected regions
\cite{where2comm,liu2020when2com} from several agents. A growing set
of methods targets what the link does to the fused result: latency
\cite{lei2022syncnet,v2xvit}, asynchrony \cite{wei2023cobevflow},
interruptions and packet loss \cite{ren2024v2xincop}, and pose error
\cite{lu2023coalign}. They are trained and evaluated against
synthetic channels: a constant delay, a Bernoulli drop, or no latency
at all. Platforms that record the link are recent and outdoor:
V2V4Real \cite{v2v4real}, V2X-Real \cite{v2xreal}, and TUMTraf-V2X
\cite{zimmer2024tumtrafv2x} provide RTK-referenced multi-agent data
with an idealized link; \cite{hawlader2024v2xexp} reports V2X
experiments on real hardware; V2X-ReaLO \cite{v2xrealo} runs online
fusion over a live Wi-Fi link; and CooperScene \cite{cooperscene}
releases per-scene C-V2X latency, throughput, and loss. CoPeD
\cite{coped}, the one indoor heterogeneous team, reports the
estimator's own pose and no link. All record the link at the
transport level, if at all, without channel state, and rely on GNSS
or on the estimator itself for pose.

\subsection{Indoor Ground Truth Without a Survey Instrument}
Indoor references differ in where the geometry comes from and whether
it is validated independently. Motion capture is accurate but
room-bound \cite{utias,m2dgr}; ultra-wideband anchors extend coverage
but need surveyed anchors and suffer in non-line-of-sight
\cite{zhao2024util}. Sparse anchoring from total-station points
\cite{hilti2022} or surveyed fiducials \cite{airmuseum} leaves
odometry to carry the pose between anchors. Map registration against
a dense prior yields full trajectories at building scale
\cite{subtmrs,lamar}, but the prior comes from a survey scanner.
SLAM-derived references dispense with the instrument and are validated
by sparse landmark distances \cite{incrowdvi} or by internal
consistency alone \cite{envodat}. SCooP-iT combines a self-built
reference map, scan-to-map registration, and map-posed fiducials, and
validates the result with measurements that never enter it
(Sec.~\ref{sec:gt}).

\subsection{Trace-Driven Network Emulation}
Record-and-replay of a measured network is established for mobile and
cellular links \cite{noble1997trace,winstein2013sprout,netravali2015mahimahi,yan2018pantheon}
and is available in the Linux kernel \cite{hemminger2005netem};
vehicular studies couple network and mobility simulators
\cite{sommer2011veins}. These tools replay a trace into a transport
stream, not into a multi-agent perception pipeline aligned to sensor
timestamps, and the traces carry no pose of the endpoints.

\subsection{Summary of the Gap}
Table~\ref{tab:gap_iotj} places the closest works on the eight
capabilities that the question of this paper requires, and
Table~\ref{tab:gap_fill} states each gap, why existing platforms leave
it open, what SCooP-iT provides, and where the paper shows it. The
wireless testbeds measure the link but not the task; the cooperative
perception platforms score the task but simulate the link or measure
it only above the transport layer, outdoors, under GNSS; the
networked-robot systems run over real links but score nothing.
SCooP-iT is the first platform in which an indoor, heterogeneous team
records per-pair link and channel state on the sensors' clock, has a
validated pose for every scan without motion capture or GNSS, and
replays the recorded link into any collaborative method.

\input{tables/gap_iotj}
\input{tables/gap_fill_iotj}

% Needs ../tables/gap_iotj.tex (capability table) and
% ../tables/gap_fill_iotj.tex (gap -> what SCooP-iT provides -> evidence).
% New bib entries are appended to ../references.bib under "IoT-J additions".
% IoT-J readers come from wireless/IoT systems, not robotics, so the order
% is: wireless testbeds -> CSI sensing -> networked robots -> cooperative
% perception -> indoor ground truth -> trace-driven emulation.
%=======================================================================
\section{Related Work}
\label{sec:related}

\subsection{Wireless and IoT Experimentation Testbeds}
Open testbeds made wireless and IoT experiments reproducible at scale.
FIT IoT-LAB \cite{adjih2015iotlab} offers thousands of low-power nodes
and over a hundred mobile robots across several indoor sites; the PAWR
platforms POWDER \cite{breen2021powder} and COSMOS
\cite{raychaudhuri2020cosmos} provide city-scale software-defined radio,
massive MIMO, and edge computing; Colosseum
\cite{bonati2021colosseum} emulates the channel in hardware for
hundreds of radios. These testbeds treat the radio as the object of
study: they measure the link precisely, but carry no perception
sensors, no per-node pose reference, and no task whose output depends
on what crosses the link. Their mobile nodes are carriers of radios,
not agents whose position and observations are scored.

\subsection{Wi-Fi Channel State Information as a Measurement}
Commodity CSI extraction, first on the Intel~5300
\cite{halperin2011csitool}, then on Broadcom chips with Nexmon
\cite{nexmon} and across several NIC families with PicoScenes
\cite{jiang2022picoscenes}, turned Wi-Fi into a sensor
\cite{ma2019wifisensing}. Channel-sounding and sensing datasets such as
DICHASUS \cite{dichasus}, ESPARGOS \cite{espargos}, MM-Fi \cite{mmfi},
and DeepSense~6G \cite{deepsense6g} pair CSI or beam power with
positions or other modalities, and ray-traced digital twins
\cite{hoydis2023sionnart,alkhateeb2023digitaltwin,bostontwin,deepmimo}
and learned radio maps \cite{levie2021radiounet} predict the channel
from geometry. In this line, CSI serves sensing or localization of a
target, and the geometry is either surveyed or simulated. None records
CSI on the links that a robot team uses to exchange perception data,
on the same clock as the robots' sensors, and none checks a measured
map against the measured channel as two independent references.

\subsection{Networked and Communication-Aware Multi-Robot Systems}
The Internet of Robotic Things \cite{simoens2018iort} treats robots as
mobile IoT nodes that sense, act, and share data over the same network
as other devices. Communication-aware robotics plans motion around the
channel \cite{muralidharan2021commaware}: robots map signal strength
online \cite{fink2010radiomap} or exploit the directionality of the
received signal to keep a team connected \cite{gil2015adaptive}.
ROS-NetSim \cite{calvofullana2021rosnetsim} couples a robot simulator
with a network simulator; CHORD and ACHORD
\cite{ginting2021chord,saboia2022achord} and MOCHA
\cite{cladera2024mocha} share data over intermittent real links in
subterranean and field deployments; \cite{blumenkamp2022framework}
provides a framework for real-robot wireless experiments; and
\cite{zhang2024ros2middleware} compares DDS, Zenoh, and MQTT for ROS~2
over Ethernet, Wi-Fi, and 4G. Multi-robot testbeds
\cite{wilson2021robotarium,paull2017duckietown,kloock2021cpmlab,testa2021choirbot}
give precise positions from motion capture or overhead cameras but
treat the network as ideal. These works model, survive, or benchmark
the link; none logs it per pair, at the rate of the sensors and on
their clock, in a form that can be replayed into another method, and
none scores perception against an independent reference.

\subsection{Cooperative Perception Under Imperfect Communication}
Cooperative perception fuses raw points \cite{chen2019cooper},
features \cite{chen2019fcooper,wang2020v2vnet}, or selected regions
\cite{where2comm,liu2020when2com} from several agents. A growing set
of methods targets what the link does to the fused result: latency
\cite{lei2022syncnet,v2xvit}, asynchrony \cite{wei2023cobevflow},
interruptions and packet loss \cite{ren2024v2xincop}, and pose error
\cite{lu2023coalign}. They are trained and evaluated against
synthetic channels: a constant delay, a Bernoulli drop, or no latency
at all. Platforms that record the link are recent and outdoor:
V2V4Real \cite{v2v4real}, V2X-Real \cite{v2xreal}, and TUMTraf-V2X
\cite{zimmer2024tumtrafv2x} provide RTK-referenced multi-agent data
with an idealized link; \cite{hawlader2024v2xexp} reports V2X
experiments on real hardware; V2X-ReaLO \cite{v2xrealo} runs online
fusion over a live Wi-Fi link; and CooperScene \cite{cooperscene}
releases per-scene C-V2X latency, throughput, and loss. CoPeD
\cite{coped}, the one indoor heterogeneous team, reports the
estimator's own pose and no link. All record the link at the
transport level, if at all, without channel state, and rely on GNSS
or on the estimator itself for pose.

\subsection{Indoor Ground Truth Without a Survey Instrument}
Indoor references differ in where the geometry comes from and whether
it is validated independently. Motion capture is accurate but
room-bound \cite{utias,m2dgr}; ultra-wideband anchors extend coverage
but need surveyed anchors and suffer in non-line-of-sight
\cite{zhao2024util}. Sparse anchoring from total-station points
\cite{hilti2022} or surveyed fiducials \cite{airmuseum} leaves
odometry to carry the pose between anchors. Map registration against
a dense prior yields full trajectories at building scale
\cite{subtmrs,lamar}, but the prior comes from a survey scanner.
SLAM-derived references dispense with the instrument and are validated
by sparse landmark distances \cite{incrowdvi} or by internal
consistency alone \cite{envodat}. SCooP-iT combines a self-built
reference map, scan-to-map registration, and map-posed fiducials, and
validates the result with measurements that never enter it
(Sec.~\ref{sec:gt}).

\subsection{Trace-Driven Network Emulation}
Record-and-replay of a measured network is established for mobile and
cellular links \cite{noble1997trace,winstein2013sprout,netravali2015mahimahi,yan2018pantheon}
and is available in the Linux kernel \cite{hemminger2005netem};
vehicular studies couple network and mobility simulators
\cite{sommer2011veins}. These tools replay a trace into a transport
stream, not into a multi-agent perception pipeline aligned to sensor
timestamps, and the traces carry no pose of the endpoints.

\subsection{Summary of the Gap}
Table~\ref{tab:gap_iotj} places the closest works on the eight
capabilities that the question of this paper requires, and
Table~\ref{tab:gap_fill} states each gap, why existing platforms leave
it open, what SCooP-iT provides, and where the paper shows it. The
wireless testbeds measure the link but not the task; the cooperative
perception platforms score the task but simulate the link or measure
it only above the transport layer, outdoors, under GNSS; the
networked-robot systems run over real links but score nothing.
SCooP-iT is the first platform in which an indoor, heterogeneous team
records per-pair link and channel state on the sensors' clock, has a
validated pose for every scan without motion capture or GNSS, and
replays the recorded link into any collaborative method.

\begin{table*}[!t]
\caption{Capabilities of Wireless Testbeds, Networked Multi-Robot Systems, and Cooperative-Perception Platforms}
\label{tab:gap_iotj}
\begin{center}
\renewcommand{\arraystretch}{1.15}
\setlength{\tabcolsep}{3pt}
\scriptsize
\begin{tabular}{@{}l cccccccc@{}}
\toprule
\textbf{Platform} & \textbf{Indoor} & \shortstack{\textbf{Mobile +}\\\textbf{infra agents}} & \shortstack{\textbf{Per-pair link}\\\textbf{state logged}} & \shortstack{\textbf{CSI on the}\\\textbf{data links}} & \shortstack{\textbf{Link and sensors}\\\textbf{on one clock}} & \shortstack{\textbf{Pose ref.\ w/o}\\\textbf{mocap/GNSS}} & \shortstack{\textbf{Multi-agent}\\\textbf{perception scoring}} & \shortstack{\textbf{Link replay}\\\textbf{into methods}} \\
\midrule
\multicolumn{9}{@{}l}{\textit{Wireless and IoT testbeds}} \\
FIT IoT-LAB \cite{adjih2015iotlab}              & \cmark & \cmark & \cmark & \xmark & \xmark & \xmark & \xmark & \xmark \\
POWDER \cite{breen2021powder}                   & \xmark & \cmark & \cmark & \cmark (SDR) & partial (GPS) & \xmark (GNSS) & \xmark & \xmark \\
COSMOS \cite{raychaudhuri2020cosmos}            & partial & \cmark & \cmark & \cmark (SDR) & partial & \xmark & \xmark & \xmark \\
Colosseum \cite{bonati2021colosseum}            & --- & --- & \cmark (emul.) & \cmark (emul.) & \xmark & --- & \xmark & \cmark (RF) \\
\midrule
\multicolumn{9}{@{}l}{\textit{Wi-Fi CSI tools and channel datasets}} \\
Nexmon / PicoScenes \cite{nexmon,jiang2022picoscenes} & --- & --- & \xmark & \cmark & \xmark & --- & \xmark & \xmark \\
DICHASUS / ESPARGOS \cite{dichasus,espargos}   & \cmark & \xmark & \xmark & \cmark & partial & \xmark (survey) & \xmark & \xmark \\
DeepSense 6G \cite{deepsense6g}                 & \xmark & partial & partial (beam) & \xmark & partial & \xmark (GPS) & \xmark & \xmark \\
\midrule
\multicolumn{9}{@{}l}{\textit{Networked and communication-aware multi-robot systems}} \\
Robotarium / CPM Lab \cite{wilson2021robotarium,kloock2021cpmlab} & \cmark & \xmark & \xmark & \xmark & \xmark & \xmark (mocap) & \xmark & \xmark \\
ROS-NetSim \cite{calvofullana2021rosnetsim}     & --- & \cmark & \xmark (sim.) & \xmark & \xmark & --- & \xmark & \cmark (sim.) \\
CHORD / ACHORD \cite{ginting2021chord,saboia2022achord} & \cmark & \cmark & partial & \xmark & \xmark & \xmark & \xmark & \xmark \\
MOCHA \cite{cladera2024mocha}                   & \cmark & \cmark & partial & \xmark & \xmark & \xmark & \xmark & \xmark \\
\midrule
\multicolumn{9}{@{}l}{\textit{Cooperative-perception platforms}} \\
CoPeD \cite{coped}                              & \cmark & \cmark & \xmark & \xmark & \xmark & \xmark (estimator) & \cmark & \xmark \\
V2V4Real \cite{v2v4real}                        & \xmark & \xmark & \xmark & \xmark & \xmark & \xmark (RTK) & \cmark & \xmark \\
TUMTraf-V2X \cite{zimmer2024tumtrafv2x}         & \xmark & \cmark & \xmark & \xmark & \xmark & \xmark (RTK) & \cmark & \xmark \\
V2X-ReaLO \cite{v2xrealo}                       & \xmark & \cmark & latency & \xmark & \xmark & \xmark (RTK) & \cmark & live \\
CooperScene \cite{cooperscene}                  & \xmark & \cmark & \cmark & \xmark & \xmark & \xmark (RTK) & \cmark & \todo{} \\
\midrule
\multicolumn{9}{@{}l}{\textit{Trace-driven network emulation}} \\
Mahimahi / NetEm \cite{netravali2015mahimahi,hemminger2005netem} & --- & --- & \cmark (trace) & \xmark & \xmark & --- & \xmark & transport \\
\midrule
\textbf{SCooP-iT (ours)}                       & \cmark & \cmark & \cmark & \cmark & \cmark & \cmark & \cmark & \cmark \\
\bottomrule
\end{tabular}
\end{center}
\vspace{2pt}
{\scriptsize \cmark: provided; partial: provided in part; \xmark: not provided; ---: not applicable. \todo{Verify every cell against the cited paper before submission, in particular the clock column for POWDER, COSMOS, and DICHASUS, and the replay cell for CooperScene.}}
\end{table*}

\begin{table*}[!t]
\caption{Gaps Left by Each Discipline and What SCooP-iT Provides}
\label{tab:gap_fill}
\begin{center}
\renewcommand{\arraystretch}{1.2}
\setlength{\tabcolsep}{4pt}
\scriptsize
\begin{tabularx}{\textwidth}{@{}>{\raggedright\arraybackslash}p{1.5cm} >{\raggedright\arraybackslash}p{2.6cm} >{\raggedright\arraybackslash}X >{\raggedright\arraybackslash}X >{\raggedright\arraybackslash}p{2.4cm}@{}}
\toprule
\textbf{Discipline} & \textbf{Gap} & \textbf{What the field has, and what it assumes away} & \textbf{What SCooP-iT provides} & \textbf{Evidence} \\
\midrule
\multirow[t]{3}{=}{\textbf{Wireless comms}}
& W1 Channel state on the links that carry the team's data
& CSI tools and datasets \cite{nexmon,jiang2022picoscenes,dichasus,espargos} sense a person or a target; V2X platforms \cite{v2xrealo,cooperscene} stop at the transport layer
& Per-pair link state at \SI{5}{\hertz}, RTT at \SI{100}{\hertz}, goodput in \SI{4}{\second} windows, and CSI at \SIrange{104}{141}{\hertz} on agent links
& Sec.~\ref{sec:link} \\
\cmidrule(l){2-5}
& W2 Measured channel checked against measured geometry
& Radio maps and digital twins predict the channel from surveyed or synthetic geometry \cite{hoydis2023sionnart,levie2021radiounet,alkhateeb2023digitaltwin}; robot-built radio maps \cite{fink2010radiomap} have no validated pose
& A measured map and a measured channel that agree although neither produced the other
& E2: Spearman $\rho = 0.60$, \SI{5.4}{\decibel} clear-to-obstructed drop \\
\cmidrule(l){2-5}
& W3 Contention and freshness on perception traffic
& 802.11 contention \cite{bianchi2000dcf} and age of information \cite{kaul2012aoi,yates2021aoi} are modeled under saturated or synthetic traffic
& Contention as a controlled knob; arrival of real perception-sized messages against a sensor-frame deadline
& E5: 86\,\% of \SI{1.5}{\mega\byte} messages within one LiDAR frame alone, \SIrange{59}{68}{\percent} with two robots \\
\midrule
\multirow[t]{3}{=}{\textbf{Robotics}}
& R1 Real links and scored perception in one experiment
& Communication-aware systems \cite{gil2015adaptive,ginting2021chord,cladera2024mocha} run over real links but score nothing against a reference; C-SLAM \cite{swarmslam,kimeramulti} is scored but its link is not logged with it
& Link logs and a pose reference for every scan of every agent on one timeline
& Secs.~\ref{sec:gt}, \ref{sec:link} \\
\cmidrule(l){2-5}
& R2 Building-scale pose reference without a survey instrument
& Motion capture is room-bound \cite{utias,m2dgr}; UWB \cite{zhao2024util} and total stations \cite{hilti2022} need surveyed anchors; outdoor platforms use RTK; CoPeD \cite{coped} reports the estimator's own pose
& Self-built map, scan-to-map registration, fiducial chain for agents without LiDAR, validated by measurements that never enter it
& \SI{1.06}{\centi\metre} surface agreement; \SI{10}{\milli\metre} (LiDAR robot), \SIrange{31}{54}{\milli\metre} (RGB-D cart) on held-out boards \\
\cmidrule(l){2-5}
& R3 Heterogeneous team with infrastructure, indoors
& Indoor teams \cite{coped,ginting2021chord,cladera2024mocha} have no infrastructure sensors; infrastructure-assisted platforms \cite{zimmer2024tumtrafv2x,v2xrealo} are outdoor
& LiDAR robot, RGB-D cart, and camera--radar infrastructure masts on one access point
& Sec.~\ref{sec:design} \\
\midrule
\multirow[t]{3}{=}{\textbf{Computing and networked systems}}
& C1 A stated bound between link logs and sensor data
& NTP \cite{mills1991ntp} and FTSP \cite{maroti2004ftsp} are evaluated on their own; no platform in Table~\ref{tab:gap_iotj} reports the offset between its network logs and its sensors
& One chrony domain for all agents and the CSI host; sniffer bias calibrated; an error budget that converts clock error into displacement at agent speed
& Measured offset $\leq$\,\SI{0.99}{\milli\second}, RMS offsets of tens of \si{\micro\second}; E1 \\
\cmidrule(l){2-5}
& C2 Middleware and edge offloading under a measured indoor channel
& ROS~2 middleware \cite{zhang2024ros2middleware} and edge perception \cite{liu2019edgeav,shao2022taskoriented,chen2019fcooper} are benchmarked with synthetic payloads or modeled channels
& ROS~2 over a contended indoor AP, with RTT, goodput, and perception-sized messages (\SI{1}{\kilo\byte} to \SI{1.5}{\mega\byte}) on the sensors' clock
& RTT 5--\SI{10}{\milli\second} idle, 25--\SI{33}{\milli\second} under load; live payload arrivals \todo{payload-node session} \\
\cmidrule(l){2-5}
& C3 Replay into a multi-agent pipeline, not a transport stream
& Trace emulators \cite{netravali2015mahimahi,hemminger2005netem,winstein2013sprout} replay into a socket; ROS-NetSim \cite{calvofullana2021rosnetsim} simulates the link
& Channel-in-the-loop replay of the recorded per-pair link into any method, keyed to sensor timestamps
& Sec.~\ref{sec:replay}; E3 \todo{replay vs live} \\
\midrule
\multirow[t]{2}{=}{\textbf{Perception}}
& P1 Evaluation against the link the agents actually had
& Methods robust to latency, asynchrony, and loss \cite{v2xvit,lei2022syncnet,wei2023cobevflow,ren2024v2xincop} are evaluated under a constant delay, a Bernoulli drop, or no latency
& Measured arrival of every message, replayable into the same methods
& E5: synthetic channels deliver 100\,\% of the messages the measured link delivers at \SIrange{59}{68}{\percent} \\
\cmidrule(l){2-5}
& P2 Real localization error on the same timeline as the link
& Pose-robust fusion \cite{lu2023coalign} is evaluated under pose noise drawn from a chosen distribution
& Each agent's estimator pose scored against the reference, on the clock of the link log
& Sec.~\ref{sec:gt} (scoring capability) \\
\midrule
\multirow[t]{2}{=}{\textbf{Testbeds}}
& T1 One platform across the four disciplines
& Wireless testbeds \cite{adjih2015iotlab,breen2021powder,raychaudhuri2020cosmos} run no perception task; multi-robot testbeds \cite{wilson2021robotarium,kloock2021cpmlab} assume an ideal network; perception platforms \cite{v2v4real,zimmer2024tumtrafv2x,cooperscene} are outdoor and GNSS-referenced
& Every column of Table~\ref{tab:gap_iotj}
& Table~\ref{tab:gap_iotj} \\
\cmidrule(l){2-5}
& T2 A link condition that reproduces, on hardware others can buy
& Testbeds report a link once; references depend on mocap, RTK, or a survey scanner
& Load mode as a named setting compared across days; commodity Wi-Fi, chrony, and printed fiducials; bill of materials
& E4 \todo{day half pending}; Table~\ref{tab:bom} \\
\bottomrule
\end{tabularx}
\end{center}
\end{table*}


% Needs ../tables/gap_iotj.tex (capability table) and
% ../tables/gap_fill_iotj.tex (gap -> what SCooP-iT provides -> evidence).
% New bib entries are appended to ../references.bib under "IoT-J additions".
% IoT-J readers come from wireless/IoT systems, not robotics, so the order
% is: wireless testbeds -> CSI sensing -> networked robots -> cooperative
% perception -> indoor ground truth -> trace-driven emulation.
%=======================================================================
\section{Related Work}
\label{sec:related}

\subsection{Wireless and IoT Experimentation Testbeds}
Open testbeds made wireless and IoT experiments reproducible at scale.
FIT IoT-LAB \cite{adjih2015iotlab} offers thousands of low-power nodes
and over a hundred mobile robots across several indoor sites; the PAWR
platforms POWDER \cite{breen2021powder} and COSMOS
\cite{raychaudhuri2020cosmos} provide city-scale software-defined radio,
massive MIMO, and edge computing; Colosseum
\cite{bonati2021colosseum} emulates the channel in hardware for
hundreds of radios. These testbeds treat the radio as the object of
study: they measure the link precisely, but carry no perception
sensors, no per-node pose reference, and no task whose output depends
on what crosses the link. Their mobile nodes are carriers of radios,
not agents whose position and observations are scored.

\subsection{Wi-Fi Channel State Information as a Measurement}
Commodity CSI extraction, first on the Intel~5300
\cite{halperin2011csitool}, then on Broadcom chips with Nexmon
\cite{nexmon} and across several NIC families with PicoScenes
\cite{jiang2022picoscenes}, turned Wi-Fi into a sensor
\cite{ma2019wifisensing}. Channel-sounding and sensing datasets such as
DICHASUS \cite{dichasus}, ESPARGOS \cite{espargos}, MM-Fi \cite{mmfi},
and DeepSense~6G \cite{deepsense6g} pair CSI or beam power with
positions or other modalities, and ray-traced digital twins
\cite{hoydis2023sionnart,alkhateeb2023digitaltwin,bostontwin,deepmimo}
and learned radio maps \cite{levie2021radiounet} predict the channel
from geometry. In this line, CSI serves sensing or localization of a
target, and the geometry is either surveyed or simulated. None records
CSI on the links that a robot team uses to exchange perception data,
on the same clock as the robots' sensors, and none checks a measured
map against the measured channel as two independent references.

\subsection{Networked and Communication-Aware Multi-Robot Systems}
The Internet of Robotic Things \cite{simoens2018iort} treats robots as
mobile IoT nodes that sense, act, and share data over the same network
as other devices. Communication-aware robotics plans motion around the
channel \cite{muralidharan2021commaware}: robots map signal strength
online \cite{fink2010radiomap} or exploit the directionality of the
received signal to keep a team connected \cite{gil2015adaptive}.
ROS-NetSim \cite{calvofullana2021rosnetsim} couples a robot simulator
with a network simulator; CHORD and ACHORD
\cite{ginting2021chord,saboia2022achord} and MOCHA
\cite{cladera2024mocha} share data over intermittent real links in
subterranean and field deployments; \cite{blumenkamp2022framework}
provides a framework for real-robot wireless experiments; and
\cite{zhang2024ros2middleware} compares DDS, Zenoh, and MQTT for ROS~2
over Ethernet, Wi-Fi, and 4G. Multi-robot testbeds
\cite{wilson2021robotarium,paull2017duckietown,kloock2021cpmlab,testa2021choirbot}
give precise positions from motion capture or overhead cameras but
treat the network as ideal. These works model, survive, or benchmark
the link; none logs it per pair, at the rate of the sensors and on
their clock, in a form that can be replayed into another method, and
none scores perception against an independent reference.

\subsection{Cooperative Perception Under Imperfect Communication}
Cooperative perception fuses raw points \cite{chen2019cooper},
features \cite{chen2019fcooper,wang2020v2vnet}, or selected regions
\cite{where2comm,liu2020when2com} from several agents. A growing set
of methods targets what the link does to the fused result: latency
\cite{lei2022syncnet,v2xvit}, asynchrony \cite{wei2023cobevflow},
interruptions and packet loss \cite{ren2024v2xincop}, and pose error
\cite{lu2023coalign}. They are trained and evaluated against
synthetic channels: a constant delay, a Bernoulli drop, or no latency
at all. Platforms that record the link are recent and outdoor:
V2V4Real \cite{v2v4real}, V2X-Real \cite{v2xreal}, and TUMTraf-V2X
\cite{zimmer2024tumtrafv2x} provide RTK-referenced multi-agent data
with an idealized link; \cite{hawlader2024v2xexp} reports V2X
experiments on real hardware; V2X-ReaLO \cite{v2xrealo} runs online
fusion over a live Wi-Fi link; and CooperScene \cite{cooperscene}
releases per-scene C-V2X latency, throughput, and loss. CoPeD
\cite{coped}, the one indoor heterogeneous team, reports the
estimator's own pose and no link. All record the link at the
transport level, if at all, without channel state, and rely on GNSS
or on the estimator itself for pose.

\subsection{Indoor Ground Truth Without a Survey Instrument}
Indoor references differ in where the geometry comes from and whether
it is validated independently. Motion capture is accurate but
room-bound \cite{utias,m2dgr}; ultra-wideband anchors extend coverage
but need surveyed anchors and suffer in non-line-of-sight
\cite{zhao2024util}. Sparse anchoring from total-station points
\cite{hilti2022} or surveyed fiducials \cite{airmuseum} leaves
odometry to carry the pose between anchors. Map registration against
a dense prior yields full trajectories at building scale
\cite{subtmrs,lamar}, but the prior comes from a survey scanner.
SLAM-derived references dispense with the instrument and are validated
by sparse landmark distances \cite{incrowdvi} or by internal
consistency alone \cite{envodat}. SCooP-iT combines a self-built
reference map, scan-to-map registration, and map-posed fiducials, and
validates the result with measurements that never enter it
(Sec.~\ref{sec:gt}).

\subsection{Trace-Driven Network Emulation}
Record-and-replay of a measured network is established for mobile and
cellular links \cite{noble1997trace,winstein2013sprout,netravali2015mahimahi,yan2018pantheon}
and is available in the Linux kernel \cite{hemminger2005netem};
vehicular studies couple network and mobility simulators
\cite{sommer2011veins}. These tools replay a trace into a transport
stream, not into a multi-agent perception pipeline aligned to sensor
timestamps, and the traces carry no pose of the endpoints.

\subsection{Summary of the Gap}
Table~\ref{tab:gap_iotj} places the closest works on the eight
capabilities that the question of this paper requires, and
Table~\ref{tab:gap_fill} states each gap, why existing platforms leave
it open, what SCooP-iT provides, and where the paper shows it. The
wireless testbeds measure the link but not the task; the cooperative
perception platforms score the task but simulate the link or measure
it only above the transport layer, outdoors, under GNSS; the
networked-robot systems run over real links but score nothing.
SCooP-iT is the first platform in which an indoor, heterogeneous team
records per-pair link and channel state on the sensors' clock, has a
validated pose for every scan without motion capture or GNSS, and
replays the recorded link into any collaborative method.

\begin{table*}[!t]
\caption{Capabilities of Wireless Testbeds, Networked Multi-Robot Systems, and Cooperative-Perception Platforms}
\label{tab:gap_iotj}
\begin{center}
\renewcommand{\arraystretch}{1.15}
\setlength{\tabcolsep}{3pt}
\scriptsize
\begin{tabular}{@{}l cccccccc@{}}
\toprule
\textbf{Platform} & \textbf{Indoor} & \shortstack{\textbf{Mobile +}\\\textbf{infra agents}} & \shortstack{\textbf{Per-pair link}\\\textbf{state logged}} & \shortstack{\textbf{CSI on the}\\\textbf{data links}} & \shortstack{\textbf{Link and sensors}\\\textbf{on one clock}} & \shortstack{\textbf{Pose ref.\ w/o}\\\textbf{mocap/GNSS}} & \shortstack{\textbf{Multi-agent}\\\textbf{perception scoring}} & \shortstack{\textbf{Link replay}\\\textbf{into methods}} \\
\midrule
\multicolumn{9}{@{}l}{\textit{Wireless and IoT testbeds}} \\
FIT IoT-LAB \cite{adjih2015iotlab}              & \cmark & \cmark & \cmark & \xmark & \xmark & \xmark & \xmark & \xmark \\
POWDER \cite{breen2021powder}                   & \xmark & \cmark & \cmark & \cmark (SDR) & partial (GPS) & \xmark (GNSS) & \xmark & \xmark \\
COSMOS \cite{raychaudhuri2020cosmos}            & partial & \cmark & \cmark & \cmark (SDR) & partial & \xmark & \xmark & \xmark \\
Colosseum \cite{bonati2021colosseum}            & --- & --- & \cmark (emul.) & \cmark (emul.) & \xmark & --- & \xmark & \cmark (RF) \\
\midrule
\multicolumn{9}{@{}l}{\textit{Wi-Fi CSI tools and channel datasets}} \\
Nexmon / PicoScenes \cite{nexmon,jiang2022picoscenes} & --- & --- & \xmark & \cmark & \xmark & --- & \xmark & \xmark \\
DICHASUS / ESPARGOS \cite{dichasus,espargos}   & \cmark & \xmark & \xmark & \cmark & partial & \xmark (survey) & \xmark & \xmark \\
DeepSense 6G \cite{deepsense6g}                 & \xmark & partial & partial (beam) & \xmark & partial & \xmark (GPS) & \xmark & \xmark \\
\midrule
\multicolumn{9}{@{}l}{\textit{Networked and communication-aware multi-robot systems}} \\
Robotarium / CPM Lab \cite{wilson2021robotarium,kloock2021cpmlab} & \cmark & \xmark & \xmark & \xmark & \xmark & \xmark (mocap) & \xmark & \xmark \\
ROS-NetSim \cite{calvofullana2021rosnetsim}     & --- & \cmark & \xmark (sim.) & \xmark & \xmark & --- & \xmark & \cmark (sim.) \\
CHORD / ACHORD \cite{ginting2021chord,saboia2022achord} & \cmark & \cmark & partial & \xmark & \xmark & \xmark & \xmark & \xmark \\
MOCHA \cite{cladera2024mocha}                   & \cmark & \cmark & partial & \xmark & \xmark & \xmark & \xmark & \xmark \\
\midrule
\multicolumn{9}{@{}l}{\textit{Cooperative-perception platforms}} \\
CoPeD \cite{coped}                              & \cmark & \cmark & \xmark & \xmark & \xmark & \xmark (estimator) & \cmark & \xmark \\
V2V4Real \cite{v2v4real}                        & \xmark & \xmark & \xmark & \xmark & \xmark & \xmark (RTK) & \cmark & \xmark \\
TUMTraf-V2X \cite{zimmer2024tumtrafv2x}         & \xmark & \cmark & \xmark & \xmark & \xmark & \xmark (RTK) & \cmark & \xmark \\
V2X-ReaLO \cite{v2xrealo}                       & \xmark & \cmark & latency & \xmark & \xmark & \xmark (RTK) & \cmark & live \\
CooperScene \cite{cooperscene}                  & \xmark & \cmark & \cmark & \xmark & \xmark & \xmark (RTK) & \cmark & \todo{} \\
\midrule
\multicolumn{9}{@{}l}{\textit{Trace-driven network emulation}} \\
Mahimahi / NetEm \cite{netravali2015mahimahi,hemminger2005netem} & --- & --- & \cmark (trace) & \xmark & \xmark & --- & \xmark & transport \\
\midrule
\textbf{SCooP-iT (ours)}                       & \cmark & \cmark & \cmark & \cmark & \cmark & \cmark & \cmark & \cmark \\
\bottomrule
\end{tabular}
\end{center}
\vspace{2pt}
{\scriptsize \cmark: provided; partial: provided in part; \xmark: not provided; ---: not applicable. \todo{Verify every cell against the cited paper before submission, in particular the clock column for POWDER, COSMOS, and DICHASUS, and the replay cell for CooperScene.}}
\end{table*}

\begin{table*}[!t]
\caption{Gaps Left by Each Discipline and What SCooP-iT Provides}
\label{tab:gap_fill}
\begin{center}
\renewcommand{\arraystretch}{1.2}
\setlength{\tabcolsep}{4pt}
\scriptsize
\begin{tabularx}{\textwidth}{@{}>{\raggedright\arraybackslash}p{1.5cm} >{\raggedright\arraybackslash}p{2.6cm} >{\raggedright\arraybackslash}X >{\raggedright\arraybackslash}X >{\raggedright\arraybackslash}p{2.4cm}@{}}
\toprule
\textbf{Discipline} & \textbf{Gap} & \textbf{What the field has, and what it assumes away} & \textbf{What SCooP-iT provides} & \textbf{Evidence} \\
\midrule
\multirow[t]{3}{=}{\textbf{Wireless comms}}
& W1 Channel state on the links that carry the team's data
& CSI tools and datasets \cite{nexmon,jiang2022picoscenes,dichasus,espargos} sense a person or a target; V2X platforms \cite{v2xrealo,cooperscene} stop at the transport layer
& Per-pair link state at \SI{5}{\hertz}, RTT at \SI{100}{\hertz}, goodput in \SI{4}{\second} windows, and CSI at \SIrange{104}{141}{\hertz} on agent links
& Sec.~\ref{sec:link} \\
\cmidrule(l){2-5}
& W2 Measured channel checked against measured geometry
& Radio maps and digital twins predict the channel from surveyed or synthetic geometry \cite{hoydis2023sionnart,levie2021radiounet,alkhateeb2023digitaltwin}; robot-built radio maps \cite{fink2010radiomap} have no validated pose
& A measured map and a measured channel that agree although neither produced the other
& E2: Spearman $\rho = 0.60$, \SI{5.4}{\decibel} clear-to-obstructed drop \\
\cmidrule(l){2-5}
& W3 Contention and freshness on perception traffic
& 802.11 contention \cite{bianchi2000dcf} and age of information \cite{kaul2012aoi,yates2021aoi} are modeled under saturated or synthetic traffic
& Contention as a controlled knob; arrival of real perception-sized messages against a sensor-frame deadline
& E5: 86\,\% of \SI{1.5}{\mega\byte} messages within one LiDAR frame alone, \SIrange{59}{68}{\percent} with two robots \\
\midrule
\multirow[t]{3}{=}{\textbf{Robotics}}
& R1 Real links and scored perception in one experiment
& Communication-aware systems \cite{gil2015adaptive,ginting2021chord,cladera2024mocha} run over real links but score nothing against a reference; C-SLAM \cite{swarmslam,kimeramulti} is scored but its link is not logged with it
& Link logs and a pose reference for every scan of every agent on one timeline
& Secs.~\ref{sec:gt}, \ref{sec:link} \\
\cmidrule(l){2-5}
& R2 Building-scale pose reference without a survey instrument
& Motion capture is room-bound \cite{utias,m2dgr}; UWB \cite{zhao2024util} and total stations \cite{hilti2022} need surveyed anchors; outdoor platforms use RTK; CoPeD \cite{coped} reports the estimator's own pose
& Self-built map, scan-to-map registration, fiducial chain for agents without LiDAR, validated by measurements that never enter it
& \SI{1.06}{\centi\metre} surface agreement; \SI{10}{\milli\metre} (LiDAR robot), \SIrange{31}{54}{\milli\metre} (RGB-D cart) on held-out boards \\
\cmidrule(l){2-5}
& R3 Heterogeneous team with infrastructure, indoors
& Indoor teams \cite{coped,ginting2021chord,cladera2024mocha} have no infrastructure sensors; infrastructure-assisted platforms \cite{zimmer2024tumtrafv2x,v2xrealo} are outdoor
& LiDAR robot, RGB-D cart, and camera--radar infrastructure masts on one access point
& Sec.~\ref{sec:design} \\
\midrule
\multirow[t]{3}{=}{\textbf{Computing and networked systems}}
& C1 A stated bound between link logs and sensor data
& NTP \cite{mills1991ntp} and FTSP \cite{maroti2004ftsp} are evaluated on their own; no platform in Table~\ref{tab:gap_iotj} reports the offset between its network logs and its sensors
& One chrony domain for all agents and the CSI host; sniffer bias calibrated; an error budget that converts clock error into displacement at agent speed
& Measured offset $\leq$\,\SI{0.99}{\milli\second}, RMS offsets of tens of \si{\micro\second}; E1 \\
\cmidrule(l){2-5}
& C2 Middleware and edge offloading under a measured indoor channel
& ROS~2 middleware \cite{zhang2024ros2middleware} and edge perception \cite{liu2019edgeav,shao2022taskoriented,chen2019fcooper} are benchmarked with synthetic payloads or modeled channels
& ROS~2 over a contended indoor AP, with RTT, goodput, and perception-sized messages (\SI{1}{\kilo\byte} to \SI{1.5}{\mega\byte}) on the sensors' clock
& RTT 5--\SI{10}{\milli\second} idle, 25--\SI{33}{\milli\second} under load; live payload arrivals \todo{payload-node session} \\
\cmidrule(l){2-5}
& C3 Replay into a multi-agent pipeline, not a transport stream
& Trace emulators \cite{netravali2015mahimahi,hemminger2005netem,winstein2013sprout} replay into a socket; ROS-NetSim \cite{calvofullana2021rosnetsim} simulates the link
& Channel-in-the-loop replay of the recorded per-pair link into any method, keyed to sensor timestamps
& Sec.~\ref{sec:replay}; E3 \todo{replay vs live} \\
\midrule
\multirow[t]{2}{=}{\textbf{Perception}}
& P1 Evaluation against the link the agents actually had
& Methods robust to latency, asynchrony, and loss \cite{v2xvit,lei2022syncnet,wei2023cobevflow,ren2024v2xincop} are evaluated under a constant delay, a Bernoulli drop, or no latency
& Measured arrival of every message, replayable into the same methods
& E5: synthetic channels deliver 100\,\% of the messages the measured link delivers at \SIrange{59}{68}{\percent} \\
\cmidrule(l){2-5}
& P2 Real localization error on the same timeline as the link
& Pose-robust fusion \cite{lu2023coalign} is evaluated under pose noise drawn from a chosen distribution
& Each agent's estimator pose scored against the reference, on the clock of the link log
& Sec.~\ref{sec:gt} (scoring capability) \\
\midrule
\multirow[t]{2}{=}{\textbf{Testbeds}}
& T1 One platform across the four disciplines
& Wireless testbeds \cite{adjih2015iotlab,breen2021powder,raychaudhuri2020cosmos} run no perception task; multi-robot testbeds \cite{wilson2021robotarium,kloock2021cpmlab} assume an ideal network; perception platforms \cite{v2v4real,zimmer2024tumtrafv2x,cooperscene} are outdoor and GNSS-referenced
& Every column of Table~\ref{tab:gap_iotj}
& Table~\ref{tab:gap_iotj} \\
\cmidrule(l){2-5}
& T2 A link condition that reproduces, on hardware others can buy
& Testbeds report a link once; references depend on mocap, RTK, or a survey scanner
& Load mode as a named setting compared across days; commodity Wi-Fi, chrony, and printed fiducials; bill of materials
& E4 \todo{day half pending}; Table~\ref{tab:bom} \\
\bottomrule
\end{tabularx}
\end{center}
\end{table*}


% Needs ../tables/gap_iotj.tex (capability table) and
% ../tables/gap_fill_iotj.tex (gap -> what SCooP-iT provides -> evidence).
% New bib entries are appended to ../references.bib under "IoT-J additions".
% IoT-J readers come from wireless/IoT systems, not robotics, so the order
% is: wireless testbeds -> CSI sensing -> networked robots -> cooperative
% perception -> indoor ground truth -> trace-driven emulation.
%=======================================================================
\section{Related Work}
\label{sec:related}

\subsection{Wireless and IoT Experimentation Testbeds}
Open testbeds made wireless and IoT experiments reproducible at scale.
FIT IoT-LAB \cite{adjih2015iotlab} offers thousands of low-power nodes
and over a hundred mobile robots across several indoor sites; the PAWR
platforms POWDER \cite{breen2021powder} and COSMOS
\cite{raychaudhuri2020cosmos} provide city-scale software-defined radio,
massive MIMO, and edge computing; Colosseum
\cite{bonati2021colosseum} emulates the channel in hardware for
hundreds of radios. These testbeds treat the radio as the object of
study: they measure the link precisely, but carry no perception
sensors, no per-node pose reference, and no task whose output depends
on what crosses the link. Their mobile nodes are carriers of radios,
not agents whose position and observations are scored.

\subsection{Wi-Fi Channel State Information as a Measurement}
Commodity CSI extraction, first on the Intel~5300
\cite{halperin2011csitool}, then on Broadcom chips with Nexmon
\cite{nexmon} and across several NIC families with PicoScenes
\cite{jiang2022picoscenes}, turned Wi-Fi into a sensor
\cite{ma2019wifisensing}. Channel-sounding and sensing datasets such as
DICHASUS \cite{dichasus}, ESPARGOS \cite{espargos}, MM-Fi \cite{mmfi},
and DeepSense~6G \cite{deepsense6g} pair CSI or beam power with
positions or other modalities, and ray-traced digital twins
\cite{hoydis2023sionnart,alkhateeb2023digitaltwin,bostontwin,deepmimo}
and learned radio maps \cite{levie2021radiounet} predict the channel
from geometry. In this line, CSI serves sensing or localization of a
target, and the geometry is either surveyed or simulated. None records
CSI on the links that a robot team uses to exchange perception data,
on the same clock as the robots' sensors, and none checks a measured
map against the measured channel as two independent references.

\subsection{Networked and Communication-Aware Multi-Robot Systems}
The Internet of Robotic Things \cite{simoens2018iort} treats robots as
mobile IoT nodes that sense, act, and share data over the same network
as other devices. Communication-aware robotics plans motion around the
channel \cite{muralidharan2021commaware}: robots map signal strength
online \cite{fink2010radiomap} or exploit the directionality of the
received signal to keep a team connected \cite{gil2015adaptive}.
ROS-NetSim \cite{calvofullana2021rosnetsim} couples a robot simulator
with a network simulator; CHORD and ACHORD
\cite{ginting2021chord,saboia2022achord} and MOCHA
\cite{cladera2024mocha} share data over intermittent real links in
subterranean and field deployments; \cite{blumenkamp2022framework}
provides a framework for real-robot wireless experiments; and
\cite{zhang2024ros2middleware} compares DDS, Zenoh, and MQTT for ROS~2
over Ethernet, Wi-Fi, and 4G. Multi-robot testbeds
\cite{wilson2021robotarium,paull2017duckietown,kloock2021cpmlab,testa2021choirbot}
give precise positions from motion capture or overhead cameras but
treat the network as ideal. These works model, survive, or benchmark
the link; none logs it per pair, at the rate of the sensors and on
their clock, in a form that can be replayed into another method, and
none scores perception against an independent reference.

\subsection{Cooperative Perception Under Imperfect Communication}
Cooperative perception fuses raw points \cite{chen2019cooper},
features \cite{chen2019fcooper,wang2020v2vnet}, or selected regions
\cite{where2comm,liu2020when2com} from several agents. A growing set
of methods targets what the link does to the fused result: latency
\cite{lei2022syncnet,v2xvit}, asynchrony \cite{wei2023cobevflow},
interruptions and packet loss \cite{ren2024v2xincop}, and pose error
\cite{lu2023coalign}. They are trained and evaluated against
synthetic channels: a constant delay, a Bernoulli drop, or no latency
at all. Platforms that record the link are recent and outdoor:
V2V4Real \cite{v2v4real}, V2X-Real \cite{v2xreal}, and TUMTraf-V2X
\cite{zimmer2024tumtrafv2x} provide RTK-referenced multi-agent data
with an idealized link; \cite{hawlader2024v2xexp} reports V2X
experiments on real hardware; V2X-ReaLO \cite{v2xrealo} runs online
fusion over a live Wi-Fi link; and CooperScene \cite{cooperscene}
releases per-scene C-V2X latency, throughput, and loss. CoPeD
\cite{coped}, the one indoor heterogeneous team, reports the
estimator's own pose and no link. All record the link at the
transport level, if at all, without channel state, and rely on GNSS
or on the estimator itself for pose.

\subsection{Indoor Ground Truth Without a Survey Instrument}
Indoor references differ in where the geometry comes from and whether
it is validated independently. Motion capture is accurate but
room-bound \cite{utias,m2dgr}; ultra-wideband anchors extend coverage
but need surveyed anchors and suffer in non-line-of-sight
\cite{zhao2024util}. Sparse anchoring from total-station points
\cite{hilti2022} or surveyed fiducials \cite{airmuseum} leaves
odometry to carry the pose between anchors. Map registration against
a dense prior yields full trajectories at building scale
\cite{subtmrs,lamar}, but the prior comes from a survey scanner.
SLAM-derived references dispense with the instrument and are validated
by sparse landmark distances \cite{incrowdvi} or by internal
consistency alone \cite{envodat}. SCooP-iT combines a self-built
reference map, scan-to-map registration, and map-posed fiducials, and
validates the result with measurements that never enter it
(Sec.~\ref{sec:gt}).

\subsection{Trace-Driven Network Emulation}
Record-and-replay of a measured network is established for mobile and
cellular links \cite{noble1997trace,winstein2013sprout,netravali2015mahimahi,yan2018pantheon}
and is available in the Linux kernel \cite{hemminger2005netem};
vehicular studies couple network and mobility simulators
\cite{sommer2011veins}. These tools replay a trace into a transport
stream, not into a multi-agent perception pipeline aligned to sensor
timestamps, and the traces carry no pose of the endpoints.

\subsection{Summary of the Gap}
Table~\ref{tab:gap_iotj} places the closest works on the eight
capabilities that the question of this paper requires, and
Table~\ref{tab:gap_fill} states each gap, why existing platforms leave
it open, what SCooP-iT provides, and where the paper shows it. The
wireless testbeds measure the link but not the task; the cooperative
perception platforms score the task but simulate the link or measure
it only above the transport layer, outdoors, under GNSS; the
networked-robot systems run over real links but score nothing.
SCooP-iT is the first platform in which an indoor, heterogeneous team
records per-pair link and channel state on the sensors' clock, has a
validated pose for every scan without motion capture or GNSS, and
replays the recorded link into any collaborative method.

\begin{table*}[!t]
\caption{Capabilities of Wireless Testbeds, Networked Multi-Robot Systems, and Cooperative-Perception Platforms}
\label{tab:gap_iotj}
\begin{center}
\renewcommand{\arraystretch}{1.15}
\setlength{\tabcolsep}{3pt}
\scriptsize
\begin{tabular}{@{}l cccccccc@{}}
\toprule
\textbf{Platform} & \textbf{Indoor} & \shortstack{\textbf{Mobile +}\\\textbf{infra agents}} & \shortstack{\textbf{Per-pair link}\\\textbf{state logged}} & \shortstack{\textbf{CSI on the}\\\textbf{data links}} & \shortstack{\textbf{Link and sensors}\\\textbf{on one clock}} & \shortstack{\textbf{Pose ref.\ w/o}\\\textbf{mocap/GNSS}} & \shortstack{\textbf{Multi-agent}\\\textbf{perception scoring}} & \shortstack{\textbf{Link replay}\\\textbf{into methods}} \\
\midrule
\multicolumn{9}{@{}l}{\textit{Wireless and IoT testbeds}} \\
FIT IoT-LAB \cite{adjih2015iotlab}              & \cmark & \cmark & \cmark & \xmark & \xmark & \xmark & \xmark & \xmark \\
POWDER \cite{breen2021powder}                   & \xmark & \cmark & \cmark & \cmark (SDR) & partial (GPS) & \xmark (GNSS) & \xmark & \xmark \\
COSMOS \cite{raychaudhuri2020cosmos}            & partial & \cmark & \cmark & \cmark (SDR) & partial & \xmark & \xmark & \xmark \\
Colosseum \cite{bonati2021colosseum}            & --- & --- & \cmark (emul.) & \cmark (emul.) & \xmark & --- & \xmark & \cmark (RF) \\
\midrule
\multicolumn{9}{@{}l}{\textit{Wi-Fi CSI tools and channel datasets}} \\
Nexmon / PicoScenes \cite{nexmon,jiang2022picoscenes} & --- & --- & \xmark & \cmark & \xmark & --- & \xmark & \xmark \\
DICHASUS / ESPARGOS \cite{dichasus,espargos}   & \cmark & \xmark & \xmark & \cmark & partial & \xmark (survey) & \xmark & \xmark \\
DeepSense 6G \cite{deepsense6g}                 & \xmark & partial & partial (beam) & \xmark & partial & \xmark (GPS) & \xmark & \xmark \\
\midrule
\multicolumn{9}{@{}l}{\textit{Networked and communication-aware multi-robot systems}} \\
Robotarium / CPM Lab \cite{wilson2021robotarium,kloock2021cpmlab} & \cmark & \xmark & \xmark & \xmark & \xmark & \xmark (mocap) & \xmark & \xmark \\
ROS-NetSim \cite{calvofullana2021rosnetsim}     & --- & \cmark & \xmark (sim.) & \xmark & \xmark & --- & \xmark & \cmark (sim.) \\
CHORD / ACHORD \cite{ginting2021chord,saboia2022achord} & \cmark & \cmark & partial & \xmark & \xmark & \xmark & \xmark & \xmark \\
MOCHA \cite{cladera2024mocha}                   & \cmark & \cmark & partial & \xmark & \xmark & \xmark & \xmark & \xmark \\
\midrule
\multicolumn{9}{@{}l}{\textit{Cooperative-perception platforms}} \\
CoPeD \cite{coped}                              & \cmark & \cmark & \xmark & \xmark & \xmark & \xmark (estimator) & \cmark & \xmark \\
V2V4Real \cite{v2v4real}                        & \xmark & \xmark & \xmark & \xmark & \xmark & \xmark (RTK) & \cmark & \xmark \\
TUMTraf-V2X \cite{zimmer2024tumtrafv2x}         & \xmark & \cmark & \xmark & \xmark & \xmark & \xmark (RTK) & \cmark & \xmark \\
V2X-ReaLO \cite{v2xrealo}                       & \xmark & \cmark & latency & \xmark & \xmark & \xmark (RTK) & \cmark & live \\
CooperScene \cite{cooperscene}                  & \xmark & \cmark & \cmark & \xmark & \xmark & \xmark (RTK) & \cmark & \todo{} \\
\midrule
\multicolumn{9}{@{}l}{\textit{Trace-driven network emulation}} \\
Mahimahi / NetEm \cite{netravali2015mahimahi,hemminger2005netem} & --- & --- & \cmark (trace) & \xmark & \xmark & --- & \xmark & transport \\
\midrule
\textbf{SCooP-iT (ours)}                       & \cmark & \cmark & \cmark & \cmark & \cmark & \cmark & \cmark & \cmark \\
\bottomrule
\end{tabular}
\end{center}
\vspace{2pt}
{\scriptsize \cmark: provided; partial: provided in part; \xmark: not provided; ---: not applicable. \todo{Verify every cell against the cited paper before submission, in particular the clock column for POWDER, COSMOS, and DICHASUS, and the replay cell for CooperScene.}}
\end{table*}

\begin{table*}[!t]
\caption{Gaps Left by Each Discipline and What SCooP-iT Provides}
\label{tab:gap_fill}
\begin{center}
\renewcommand{\arraystretch}{1.2}
\setlength{\tabcolsep}{4pt}
\scriptsize
\begin{tabularx}{\textwidth}{@{}>{\raggedright\arraybackslash}p{1.5cm} >{\raggedright\arraybackslash}p{2.6cm} >{\raggedright\arraybackslash}X >{\raggedright\arraybackslash}X >{\raggedright\arraybackslash}p{2.4cm}@{}}
\toprule
\textbf{Discipline} & \textbf{Gap} & \textbf{What the field has, and what it assumes away} & \textbf{What SCooP-iT provides} & \textbf{Evidence} \\
\midrule
\multirow[t]{3}{=}{\textbf{Wireless comms}}
& W1 Channel state on the links that carry the team's data
& CSI tools and datasets \cite{nexmon,jiang2022picoscenes,dichasus,espargos} sense a person or a target; V2X platforms \cite{v2xrealo,cooperscene} stop at the transport layer
& Per-pair link state at \SI{5}{\hertz}, RTT at \SI{100}{\hertz}, goodput in \SI{4}{\second} windows, and CSI at \SIrange{104}{141}{\hertz} on agent links
& Sec.~\ref{sec:link} \\
\cmidrule(l){2-5}
& W2 Measured channel checked against measured geometry
& Radio maps and digital twins predict the channel from surveyed or synthetic geometry \cite{hoydis2023sionnart,levie2021radiounet,alkhateeb2023digitaltwin}; robot-built radio maps \cite{fink2010radiomap} have no validated pose
& A measured map and a measured channel that agree although neither produced the other
& E2: Spearman $\rho = 0.60$, \SI{5.4}{\decibel} clear-to-obstructed drop \\
\cmidrule(l){2-5}
& W3 Contention and freshness on perception traffic
& 802.11 contention \cite{bianchi2000dcf} and age of information \cite{kaul2012aoi,yates2021aoi} are modeled under saturated or synthetic traffic
& Contention as a controlled knob; arrival of real perception-sized messages against a sensor-frame deadline
& E5: 86\,\% of \SI{1.5}{\mega\byte} messages within one LiDAR frame alone, \SIrange{59}{68}{\percent} with two robots \\
\midrule
\multirow[t]{3}{=}{\textbf{Robotics}}
& R1 Real links and scored perception in one experiment
& Communication-aware systems \cite{gil2015adaptive,ginting2021chord,cladera2024mocha} run over real links but score nothing against a reference; C-SLAM \cite{swarmslam,kimeramulti} is scored but its link is not logged with it
& Link logs and a pose reference for every scan of every agent on one timeline
& Secs.~\ref{sec:gt}, \ref{sec:link} \\
\cmidrule(l){2-5}
& R2 Building-scale pose reference without a survey instrument
& Motion capture is room-bound \cite{utias,m2dgr}; UWB \cite{zhao2024util} and total stations \cite{hilti2022} need surveyed anchors; outdoor platforms use RTK; CoPeD \cite{coped} reports the estimator's own pose
& Self-built map, scan-to-map registration, fiducial chain for agents without LiDAR, validated by measurements that never enter it
& \SI{1.06}{\centi\metre} surface agreement; \SI{10}{\milli\metre} (LiDAR robot), \SIrange{31}{54}{\milli\metre} (RGB-D cart) on held-out boards \\
\cmidrule(l){2-5}
& R3 Heterogeneous team with infrastructure, indoors
& Indoor teams \cite{coped,ginting2021chord,cladera2024mocha} have no infrastructure sensors; infrastructure-assisted platforms \cite{zimmer2024tumtrafv2x,v2xrealo} are outdoor
& LiDAR robot, RGB-D cart, and camera--radar infrastructure masts on one access point
& Sec.~\ref{sec:design} \\
\midrule
\multirow[t]{3}{=}{\textbf{Computing and networked systems}}
& C1 A stated bound between link logs and sensor data
& NTP \cite{mills1991ntp} and FTSP \cite{maroti2004ftsp} are evaluated on their own; no platform in Table~\ref{tab:gap_iotj} reports the offset between its network logs and its sensors
& One chrony domain for all agents and the CSI host; sniffer bias calibrated; an error budget that converts clock error into displacement at agent speed
& Measured offset $\leq$\,\SI{0.99}{\milli\second}, RMS offsets of tens of \si{\micro\second}; E1 \\
\cmidrule(l){2-5}
& C2 Middleware and edge offloading under a measured indoor channel
& ROS~2 middleware \cite{zhang2024ros2middleware} and edge perception \cite{liu2019edgeav,shao2022taskoriented,chen2019fcooper} are benchmarked with synthetic payloads or modeled channels
& ROS~2 over a contended indoor AP, with RTT, goodput, and perception-sized messages (\SI{1}{\kilo\byte} to \SI{1.5}{\mega\byte}) on the sensors' clock
& RTT 5--\SI{10}{\milli\second} idle, 25--\SI{33}{\milli\second} under load; live payload arrivals \todo{payload-node session} \\
\cmidrule(l){2-5}
& C3 Replay into a multi-agent pipeline, not a transport stream
& Trace emulators \cite{netravali2015mahimahi,hemminger2005netem,winstein2013sprout} replay into a socket; ROS-NetSim \cite{calvofullana2021rosnetsim} simulates the link
& Channel-in-the-loop replay of the recorded per-pair link into any method, keyed to sensor timestamps
& Sec.~\ref{sec:replay}; E3 \todo{replay vs live} \\
\midrule
\multirow[t]{2}{=}{\textbf{Perception}}
& P1 Evaluation against the link the agents actually had
& Methods robust to latency, asynchrony, and loss \cite{v2xvit,lei2022syncnet,wei2023cobevflow,ren2024v2xincop} are evaluated under a constant delay, a Bernoulli drop, or no latency
& Measured arrival of every message, replayable into the same methods
& E5: synthetic channels deliver 100\,\% of the messages the measured link delivers at \SIrange{59}{68}{\percent} \\
\cmidrule(l){2-5}
& P2 Real localization error on the same timeline as the link
& Pose-robust fusion \cite{lu2023coalign} is evaluated under pose noise drawn from a chosen distribution
& Each agent's estimator pose scored against the reference, on the clock of the link log
& Sec.~\ref{sec:gt} (scoring capability) \\
\midrule
\multirow[t]{2}{=}{\textbf{Testbeds}}
& T1 One platform across the four disciplines
& Wireless testbeds \cite{adjih2015iotlab,breen2021powder,raychaudhuri2020cosmos} run no perception task; multi-robot testbeds \cite{wilson2021robotarium,kloock2021cpmlab} assume an ideal network; perception platforms \cite{v2v4real,zimmer2024tumtrafv2x,cooperscene} are outdoor and GNSS-referenced
& Every column of Table~\ref{tab:gap_iotj}
& Table~\ref{tab:gap_iotj} \\
\cmidrule(l){2-5}
& T2 A link condition that reproduces, on hardware others can buy
& Testbeds report a link once; references depend on mocap, RTK, or a survey scanner
& Load mode as a named setting compared across days; commodity Wi-Fi, chrony, and printed fiducials; bill of materials
& E4 \todo{day half pending}; Table~\ref{tab:bom} \\
\bottomrule
\end{tabularx}
\end{center}
\end{table*}

